Trees\textemdash or, more generally inductively-defined structures\textemdash are core to the study of programming languages.
This is, in part, why functional programming languages like OCaml are well-suited for \textit{implementing} programming languages.
As such, we have to spend some time on the humble notion of trees.
This appendix covers a small slice of the topic, only what we need for the main part of the text.\footnote{
There are whole books dedicated to the study of mathematical induction and inductively defined structures.
}

If you've taken a course in discrete mathematics, you've likely seen the graph-theoretic definition of trees.

\begin{definition}
  A \textbf{tree} is undirected graph which is connected and acyclic.
  A \textbf{directed tree} is a directed graph whose underlying graph is a tree.
  A directed tree is \textbf{rooted} if there is a unique vertex with in-degree 0.
\end{definition}
%
We'll be primarily interested in \textit{nonempty rooted directed} trees.\footnote{
Moving forward \enquote{tree} we will always mean \enquote{nonempty rooted directed tree.}
}
These are also sometimes called \textbf{rose trees}.
To make this more explicit we'll work with the \textit{inductive} definition of (rose) trees.

\begin{definition}
  \label{def:rose-tree}
  A \textbf{tree} with values from $V$ is defined inductively as follows:
  \begin{itemize}
  \item if $v$ is a value from $V$ and $T_1, \dots, T_k$ are trees then so is $\tnode(v, T_1, \dots, T_k)$.
  \end{itemize}
  Note, in particular, that $\tnode(v)$ a tree for any value $v$ from $V$.
  We call this kind of tree a \textbf{leaf}.
  The \textbf{root} of a tree is defined as:
  \begin{align*}
    \troot {\tnode(v, T_1, \dots, T_k)} \defeq v
  \end{align*}
  and the \textbf{children} a tree are defined as:
  \begin{align*}
    \tchildren{\tnode(v, T_1, \dots, T_k)} \defeq (T_1, \dots, T_k)
  \end{align*}
\end{definition}

\begin{example}
  \label{ex:decision-tree}
  Decision trees represent boolean functions, and we can represent a decisions tree as a rose tree.
  One decision tree for the boolean function $\mathsf{OR}(x_1, x_2, x_3) = x_1 \lor x_2 \lor x_3$ is:
  \begin{align*}
    \tnode(x_1 = 1, \tnode(1), \tnode(x_2 = 1, \tnode(1), \tnode(x_3 = 1, \tnode(1), \tnode(0))))
  \end{align*}
  Values of this tree are queries to the inputs of the function.
  Roughly speaking, this decision tree expresses that the value of $\mathsf{OR}(x_1, x_2, x_3)$ can be determined by searching from left to right for an input which is equal to 1.
\end{example}
%
We can naturally define rose trees in OCaml, which means we can also define the usual functions on trees:
\begin{ocaml}
  type 'a rose_tree = Node of 'a * 'a rose_tree list
  let example = Node (1, [Node (2, []); Node (3, [])])

  let rec depth (Node (_, ts)) =
    List.fold_left
      (fun acc n -> max acc (n + 1))
      0
      (List.map depth ts)

  let rec size (Node (_, ts)) =
    List.(fold_left (+) 1 (map size ts))

  let _ = assert (depth example = 1)
  let _ = assert (size example = 3)
\end{ocaml}
%
We visualize trees in the usual way.
Here, for example, is a visualization of the tree from \Cref{ex:decision-tree}:

\begin{center}
  \begin{forest}
    for tree={
      minimum width=1in,
    }
    [{$x_1 = 1$}
      [$1$]
      [{$x_2 = 1$}
        [$1$]
        [{$x_3 = 1$}
          [$1$]
          [$0$]
        ]
      ]
    ]
  \end{forest}
\end{center}
%
We'll also occassionally use the following form of tree visualization, which we'll call \textbf{compact form}:

\begin{center}
  \begin{forest}
    for tree={
      grow'=0,
      child anchor=west,
      parent anchor=south,
      anchor=west,
      calign=first,
      edge path={
        \noexpand\path [draw, \forestoption{edge}]
        (!u.south west) +(5pt,0) |- (.child anchor)\forestoption{edge label};
      },
      before typesetting nodes={
        if n=1
        {insert before={[,phantom]}}
        {}
      },
      before computing xy={l=20pt},
    }
    [{$x_1 = 1$}
      [$1$]
      [{$x_2 = 1$}
        [$1$]
        [{$x_3 = 1$}
          [$1$]
          [$0$]
        ]
      ]
    ]
  \end{forest}
\end{center}
%
This form is based on visualizations of file trees.\footnote{
If you're interested, look up the command line tool \texttt{tree}, or just type \texttt{tree .} in your terminal and see what you get.
}
It's particularly useful when a tree it too wide to draw otherwise, e.g., when the values in nodes are themselves very wide.
There's one more form of tree visualization we'll use for drawing derivation trees; this will be covered in \Cref{ch:infer}.

\section{Structural Induction}


Induction is a principle used to prove universal statements about inductively-defined structures (like trees).
If you've taken a course in discrete mathematics, you've likely seen the principle of natural number induction:

\begin{quote}
  \textit{
    Given a property $P$ of natural numbers, if:
    \begin{itemize}
    \item
      $P$ holds of the number $0$;
    \item
      for every number $k$, if we assume $P$ holds of $k$ (an assumption often called the \textbf{induction hypothesis}), then we can demonstrate that $P$ holds of $k + 1$;
    \end{itemize}
    then $P$ holds of every natural number.
  }
\end{quote}

\begin{exercise}
  Prove that
  \begin{align*}
    2\sum_{i = 1}^n i = n(n+1)
  \end{align*}
  holds for all natural numbers $n$ using natural number induction.
\end{exercise}
%
Inductive structures\textemdash by definition of being inductive\textemdash have analogous induction principles.
Here's the principle of tree induction:

\begin{quote}
  \textit{
    Given a property $P$ of trees with values from $V$, if
    \begin{itemize}
    \item for any value $v$ from $V$, and trees $T_1, \dots, T_k$, if we assume that $P$ holds of each tree $T_1, \dots, T_k$\footnote{Again, this assumption is called the \textit{induction hypothesis}.} then we can demonstrate that $P$ holds of $\tnode(v, T_1, \dots, T_k)$;
    \end{itemize}
    then $P$ holds of all trees with values from $V$.
  }
\end{quote}
%
We won't concern ourselves further with the general notion of trees.\footnote{There isn't much more we can say without also formally defining \textit{recursion} on trees, another interesting topic that we'll unceremoniously brush under the rug.}
We'll focus on \textbf{inductively-define subsets} of trees.

\begin{notation}
  For a set $A$, we write $A^*$ for the set of finite sequences of elements of $A$ and we write $A^+$ for the subset of nonempty sequences in $A^*$.
  For example, the sequence $(1, 2, 3, 4, 5)$ is an element of both $\N^*$ and $\N^+$.
\end{notation}

\begin{definition}
  Let $Q$ be a property on $V^+$.
  The \textbf{inductively-defined subset} of trees $\mathcal Q$ given by $Q$ is defined (inductively) as follows:
  \begin{itemize}
  \item for any value $v$ from $V$ and trees $T_1, \dots, T_k$ from $\mathcal Q$, if $Q$ holds of $(v, \troot {v_1}, \dots, \troot {v_k})$ then $\tnode(v, T_1, \dots, T_k) \in \mathcal Q$.
  \end{itemize}
\end{definition}
%
In other words, inductively-defined subsets of trees are defined by describing when we're allowed to build new trees from those we've already built.

\begin{example}
  \label{ex:prime-fact}
  Consider the inductively-defined set $\mathcal Q$ of trees with values from $\N$ given by the property $Q$ which holds of $(v, v_1,\dots, v_k)$ when $v$ is prime or $v = v_1 \dots v_k$.
  That is, the value at the root of a given tree is prime or the product of the roots of its children.
  The following is an example of such a tree.
  \begin{center}
    \begin{forest}
      [$60$
        [$5$]
        [$12$
          [$2$]
          [$3$]
          [$2$]
        ]
      ]
    \end{forest}
  \end{center}
  We might recognize that the leaves of this tree define a prime factorization of the root, e.g., $60 = 5 * 2 * 3 * 2$.
\end{example}
%
Once we have inductively-defined subsets of trees, we can define their corresponding principle of induction.
It's the same as the principle of tree induction but with a slightly stronger induction hypothesis:

\begin{quote}
  \textit{
    Given a property $P$ of trees with values from $V$ \textcolor{red}{and a property $Q$ which inductively-defines a subset $\mathcal Q$ of trees}, if
    \begin{itemize}
    \item for any value $v$ from $V$, and trees $T_1, \dots, T_k$ \textcolor{red}{from $\mathcal Q$}, if we assume that $P$ holds of each tree $T_1, \dots, T_k$ \textcolor{red}{and that $Q$ holds of the sequence $(v, \troot {T_1}, \dots, \troot {T_k})$} then we can demonstrate that $P$ holds of $\tnode(v, T_1, \dots, T_k)$;
    \end{itemize}
    then $P$ holds of all trees \textcolor{red}{in $\mathcal Q$}.
  }
\end{quote}
%
Let's see this in action.
Consider the subset $\mathcal Q$ of trees defined in \Cref{ex:prime-fact}.
We'd like to prove that for any tree $T$ from $\mathcal Q$, the number $\troot{T}$ has a prime factorization.

\begin{proof}
Let $\tnode(n, T_1, \dots, T_k)$ be a tree in $\mathcal Q$ and suppose that $\troot {T_i}$ has a prime factorization $p_{1, i}\dots p_{l_i, i}$ for all indices $i$.
Furthermore we assume, by definition of $\mathcal Q$, that $n$ is prime or
\begin{align*}
  n = \troot{T_1}\dots\troot{T_k}
\end{align*}
If $n$ is prime, then $n$ is also its prime factorization.
Otherwise,
\begin{align*}
  n = \prod_{i = 1}^k\troot{T_i} = \prod_{i = 1}^k (p_{1, i}\dots p_{l_i, i})
\end{align*}
yielding a prime factorization of $n$.
That is, we take the product of all the prime factorizations of the children of $T$ to get a prime factorization of its root.\footnote{
What we've done here is only \textit{part of} the fundamental theorem of arithmetic.
We would also need to prove that every natural number is the root of some tree in $\mathcal Q$, which we can prove by (strong) induction over natural numbers (we also say nothing of uniqueness...).
}
\end{proof}

\section{Induction on Derivations}

What we've done so far is more general than necessary, and is ultimately in service of defining the \textit{derivation induction principle}.
First, observe that the set of derivations in a given inference system $\mathcal I$, as defined in \Cref{ch:infer}, is an inductively-define subset of trees.
The property is straightforward: it holds of the nonempty sequence of judgments $(J, J_1, \dots, J_k)$ when
\begin{prooftree}
  \AxiomC{$J_1$}
  \AxiomC{$\dots$}
  \AxiomC{$J_k$}
  \TrinaryInfC{$J$}
\end{prooftree}
is (an instance of) an inference rule of $\mathcal I$.
In other words, a derivation is a tree in which every node is either an axiom or follows from an inference rule.
The upshot is that the derivation induction principle is a specialization of the induction principle for inductively-defined subsets of trees.

Our second observation: being able to prove that a property holds of all derivations is enough to be able to prove that a property holds of all \textit{derivable judgments}.
There's is an obvious correspondence between derivations and derivable judgments: if a judgment is derivable, then it has a derivation.
And this is ultimately what we care about when we reason about things like type safety (\Cref{ch:types}).
All said, we can state the derivation induction principle as follows:

\begin{quote}
  \textit{
    Given a property $P$ of judgments of $\mathcal I$, if
    \begin{itemize}
    \item for any judgment $J$ and derivable judgments $J_1, \dots, J_k$, if we assume that $P$ holds of each judgment $J_1, \dots, J_k$ where
      \begin{prooftree}
        \AxiomC{$J_1$}
        \AxiomC{$\dots$}
        \AxiomC{$J_k$}
        \TrinaryInfC{$J$}
      \end{prooftree}
    then we can demonstrate that $P$ holds of $J$;
    \end{itemize}
    then $P$ holds of all derivable judgments of $\mathcal D$.
  }
\end{quote}
%
This principle tells us that we can look at the \textit{last} inference rule applied and prove that the property holds assuming it holds of the \textit{antecedents} of the applied rule.

It may be easier to grok this principle by example.
Consider the following basic inference system over judgments of the form \enquote{$n$ is even} where $n \in \N$.

\begin{example}
  This system is taken directly from \Cref{ch:infer}.
  \begin{center}
    \AxiomC{}
    \RightLabel{\small (zero)}
    \UnaryInfC{$0$ is even}
    \DisplayProof
    \qquad
    \AxiomC{$m$ is even}
    \AxiomC{\hl{$n = m + 2$}}
    \RightLabel{\small (addTwo)}
    \BinaryInfC{$n$ is even}
    \DisplayProof
  \end{center}
\end{example}
We'd like to prove that this inference system is sound, i.e., that if \enquote{$n$ is even} is derivable, then $n$ is, in fact, even (in particular, \enquote{$3$ is even} is not derivable).
We can prove this by induction on derivations.
\begin{proof}
  Let $J$ denote an arbitrary judgment \enquote{$n$ is even.}
  There are two cases to consider.
  \begin{itemize}
  \item If $J$ follows from (zero), then $J$ must be the judgment \enquote{$0$ is even} in which case the property holds.
  \item
    If $J$ follows from (addTwo), then \enquote{$n - 2$ is even} must be derivable and we may assume that $n - 2$ is even.
    This implies that $n$ itself is even.
  \end{itemize}
\end{proof}
%
In a sense, this example might be \textit{too} simple.
It's not immediately clear that we've done anything in this proof; the distinction between $n$ being even and the judgment \enquote{$n$ is even} being derivable is admittedly a subtle one.
But, in order to avoid ballooning this appendix into a full-blown chapter, I'll relegate the presentation of more interesting examples to \Cref{ch:types}.
